% !TeX program = xelatex
\documentclass[11pt,a4paper,fontset=windows]{ctexart}

\usepackage{geometry}
\geometry{left=2.7cm,right=2.7cm,top=2.7cm,bottom=2.7cm}
\usepackage{amsmath,amssymb}
\usepackage{booktabs}
\usepackage{graphicx}
\usepackage{array}
\usepackage{caption}
\usepackage{xcolor}
\usepackage[hidelinks]{hyperref}
\graphicspath{{../outputs/}{outputs/}}

\captionsetup{font=small,labelsep=quad}

\newcommand{\R}{R^2_{\mathrm{oos}}}
\newcommand{\code}[1]{\texttt{#1}}
\newcommand{\repourl}{https://github.com/GeneralSunz/gkx2020-nn3-replication}

\title{\bfseries 三层前馈神经网络的股票收益预测\\[2pt]
\large ——Gu, Kelly \& Xiu (2020) 年度样本外 $R^2$ 的复现与可复现性边界}
\author{金融机器学习 · 期末项目\\[6pt]
\begin{tabular}{ll}
小组成员： & \underline{\hspace{4cm}} \\[2pt]
  & \underline{\hspace{4cm}} \\[2pt]
  & \underline{\hspace{4cm}} \\
\end{tabular}\\[10pt]
\normalsize 代码与复现材料：\url{\repourl}}
\date{2026 年 11 月}

\begin{document}
\maketitle

\begin{abstract}
\sloppy
\noindent
本文复现 Gu, Kelly 与 Xiu（2020，\emph{RFS}）表 2 中三层前馈神经网络（NN3）
在年度预测时域下的样本外 $R^2$，并回答一个复现工作中最关键的问题：
\textbf{当数值达不到原文时，应当归因于实现缺陷，还是归因于数据边界？}

结论分三层。第一，\textbf{月度口径可以复现}：本文 NN3 取得 $+0.47\%$，
论文为 $0.40\%$，独立第三方预注册复现为 $0.454\%$，三者互相印证，
说明从数据读取到 $R^2$ 计算的全链路正确。
第二，\textbf{年度口径未能复现}：本文最优结果为缩尾口径 $+2.71\%$、原始口径 $+1.61\%$，
论文为 $+3.40\%$。
第三，\textbf{差距不是实现缺陷}，有三项独立证据：
合成信号注入测试显示本管线可恢复已知信号理论上限的 $99\%$；
由月度信号推出的年度可预测性上界约为 $+0.84\%$，而本文年度模型已达 $+1.32\%$，
已超出该上界；论文数值要求横截面信息系数达到 $0.185$，约为月度信号可持续支撑水平的两倍。

本文进一步调查了六个公开复现项目，发现它们全部针对月度口径，
尚无公开工作复现表 2 的年度数值；其中最严谨的一项亦未能复现论文的线性基准，
与本文遇到的困难同源。据此本文认为，年度数值的未达反映的是公开数据条件下该指标的
可复现性边界，并给出了该边界的定量刻画。

\vspace{6pt}
\noindent\textbf{关键词：} 实证资产定价；机器学习；神经网络；样本外 $R^2$；可复现性
\end{abstract}

\section{引言}

Gu, Kelly 与 Xiu（2020，以下简称 GKX）系统比较了十三种机器学习方法在美股横截面收益
预测上的表现，核心结论有三：非线性方法显著优于线性基准；大市值股票更可预测；
机器学习预测具有可观的经济价值。该文公开了构造好的 94 个公司特征面板，使复现成为可能。

本文选择其中作者称为``头条模型''的三层前馈神经网络（NN3），
目标复现其在年度预测时域下的样本外 $R^2$（原文表 2）。
选择年度而非月度，是因为原文称年度 $R^2$``比月度大近一个数量级''，信噪比更高，
更容易检验模型是否学到了真实的横截面结构。

复现工作最容易犯的错误，是把``数值不等于原文''直接归因于实现缺陷；
同样有害的是反过来，用``数据不行''掩盖真正的代码错误——
后者在本项目中确实发生过：批量归一化的位置曾被写反，早停准则曾完全失效。
因此本文的组织方式不是``先给结果再解释差距''，而是先建立证据链，
把三个问题分开回答：
\begin{enumerate}
  \item 数据是否正确？（第 \ref{sec:data} 节，五项独立验证）
  \item 代码是否正确？（第 \ref{sec:results} 节，月度交叉验证 + 合成信号注入）
  \item 差距能否归因？（第 \ref{sec:results} 节、第 \ref{sec:lit} 节，定量上界 + 文献对照）
\end{enumerate}

\textbf{本文的主要发现是：月度口径可以复现且与第三方高度一致；年度口径无法复现，
但三项独立证据表明这不是实现缺陷。}公开文献中亦无人报告过年度数值。

\section{模型}\label{sec:model}

NN3 由三个隐藏层构成，神经元数依次为 $32$、$16$、$8$，呈逐层收缩的金字塔结构：
\begin{equation}
h^{(\ell)} = \mathrm{BN}\!\left(\mathrm{ReLU}(W_\ell h^{(\ell-1)} + b_\ell)\right),\qquad
\hat y = W_4 h^{(3)} + b_4 ,
\end{equation}
输出层为线性。

此处有一个容易写反的实现细节：原文 Internet Appendix 算法 7 明确指出，
批量归一化施加于 \emph{ReLU 之后的每个激活值}上，即``线性 $\to$ ReLU $\to$ BN''，
而非许多深度学习框架的常见顺序``线性 $\to$ BN $\to$ ReLU''。
本文初期采用了后者，核对附录原文后改正。该改动还带来约 $12\%$ 的训练加速。

目标函数为带 $L_1$ 惩罚的均方误差，惩罚项施加于全部网络权重；
优化器为 Adam，采用其默认超参数；训练按批大小 $10{,}000$ 进行小批量随机抽样。
$L_1$ 惩罚的作用是把不重要的连接压向零，在本问题这种信噪比极低（$R^2$ 仅为百分之一量级）
的场景中，正则化强度对结果的影响远大于网络宽度。

\section{数据与预处理}\label{sec:data}

\subsection{数据来源}

\textbf{公司特征}取自 GKX 公开的 \code{datashare.csv}：94 个公司特征、股票代码与两位 SIC
行业代码，覆盖 1957 年 3 月至 2016 年 12 月共 718 个月、3{,}760{,}208 条观测、
29{,}892 只股票。该文件未被预先过滤——平均每条观测有 $58\%$ 的特征缺失，
其中 $62\%$ 的观测缺失过半，与``按可计算特征筛选''的样本相符。

\textbf{股票收益}取自课程提供的 CRSP 月度提取文件，含 36{,}408 只股票、1925 年 12 月至
2022 年 3 月的调整后收益与市值。该文件的月份字段 \code{md} 不是日历月份，
而是授课教师预处理得到的连续月序号，其生成规则为
\begin{equation}
\code{md} = \mathrm{year}(\text{date})\times 12 + \mathrm{month}(\text{date}) - 1925\times 12 - 11,
\end{equation}
即 \code{md} $=1$ 对应 1925 年 12 月。将其反解可得
\begin{equation}
\mathrm{year}\times 12 + \mathrm{month} = (1925\times 12 + 11) + \code{md} - 1 ,
\end{equation}
本文的还原即按此式实现。校验结果：文件中 1{,}156 个不同的 \code{md} 取值
与教师的公式逐一对应，无一分歧；此外，以 \code{mom1m} 反推收益做偏移扫描，
偏移为 0 时相关 $0.919$（其余偏移约为 0），两种方法互相印证。

\textbf{无风险利率}取自 Fama--French 数据图书馆的月度因子文件
（\code{F-F\_Research\_Data\_Factors}），用于构造超额收益。
校验结果：本文所用序列与官方文件的 \code{RF} 列在 718 个重叠月份上逐月一致，
最大绝对差为 $1.0\times10^{-16}$（浮点精度极限），\code{MktRF} 亦如此。

\subsection{特征构造}

对每个月、每个特征做横截面排序并映射到 $[-1,1]$：
\begin{equation}
c_{i,t} = 2\cdot\frac{\mathrm{rank}(x_{i,t})}{N_t+1} - 1 ,
\end{equation}
缺失值在排序前用当月横截面中位数填补，填补后位于分布中央，不引入方向性偏差。
原文的完整设计为 $920$ 维（94 特征 + 94$\times$8 宏观交互 + 74 行业虚拟变量）。

本文主结果使用\textbf{仅 94 个原始特征}。理由有二：其一，宏观交互项在本框架下并未提升
验证集表现（920 维的验证集信息系数仅 $0.015$--$0.042$，远低于 94 维的 $0.11$--$0.13$，
见第 \ref{sec:gap} 节）；其二，课程要求以 94 个公司特征为主。该偏离在附录 \ref{app:deviation} 列明。

\subsection{目标变量}

年度目标为未来 12 个月的复利超额收益：
\begin{equation}
y_{i,t}^{(12)} = \prod_{k=0}^{11}\left(1+r_{i,t+k}\right) - \prod_{k=0}^{11}\left(1+r^f_{t+k}\right),
\end{equation}
窗口起点为第 $t$ 月，因为特征在第 $t$ 月末才可观测。月度目标为 $y_{i,t}^{(1)} = r_{i,t} - r^f_t$，
用于与论文表 1 及第三方复现交叉验证。

\subsection{样本与数据验证}\label{sec:validation}

样本为 3{,}737{,}632 条股票—月度观测、29{,}845 只股票，月均 5{,}206 只。
需要说明的是，原文称月均股票数``超过 6{,}200''，而作者公开的特征文件本身月均为 5{,}237，
教师提供的完整 CRSP 文件月均为 5{,}601，三者并不一致。
可以确认的是，\code{datashare.csv} 中的 29{,}886 个股票代码 $100\%$ 包含于 CRSP 文件，
两文件是包含关系而非互补关系，因此 5{,}206 是公开数据所能支持的最大估计样本；
原文的 6{,}200 计的是收益样本而非估计样本。

在建模之前，本文对数据链路做了五项独立验证（表 \ref{tab:audit}）。
这类错误往往不会以程序崩溃的形式暴露，而会安静地污染每一个下游结果。

\begin{table}[htbp]
\centering
\caption{数据链路的独立验证}
\label{tab:audit}
\small
\begin{tabular}{p{3.2cm}p{5.6cm}p{5.2cm}}
\toprule
检查项 & 方法 & 结果 \\
\midrule
月度收益分布 & 与 CRSP 已知统计性质对照 &
均值 $+1.108\%$、标准差 $17.29\%$、取值 $[-1,+24]$，与文献值一致 \\
特征排序 & 从原始 CSV 独立重算某月排名 &
与缓存设计矩阵最大差 $3\times10^{-4}$（单精度极限） \\
收益与特征对齐 & \code{mom1m} 与调整后收益的偏移扫描 &
偏移为 0 时相关 $0.919$，其余偏移约为 0 \\
目标窗口起点 & 用三个已知有效特征的横截面信息系数判定 &
当月起点 $0.0542$ 优于错后一月 $0.0455$ \\
年度目标复利 & 手工计算 12 个月连乘并比对 &
最大差 $3\times10^{-7}$ \\
月份字段 \code{md} & 按教师给定公式反解并与本实现对拍 &
1{,}156 个月份零分歧 \\
无风险利率 & 与 Fama--French 官方文件逐月比对 &
718 个月最大差 $1.0\times10^{-16}$ \\
\bottomrule
\end{tabular}
\end{table}

其中``目标窗口起点''一项值得展开。若特征实际包含当月收益，
以第 $t$ 月为窗口起点就会引入前视偏差，这是资产定价复现中最常见的错误之一。
本文用 \code{mom12m}、\code{mvel1}、\code{retvol} 三个已知有预测力的特征，
分别计算它们与两种窗口定义下年度收益的横截面秩相关，以第 $t$ 月起算者在三个特征上均更高，
说明特征确实在第 $t$ 月末可观测，不构成前视。

\section{训练与检验设计}\label{sec:design}

\subsection{划分与重估}

严格按原文的三段式时间划分：训练 1957--1974（18 年）、验证 1975--1986（12 年）、
测试 1987--2016（30 年）。训练窗口逐年扩张，验证窗口固定 12 年并随之滚动，
模型每年重估一次。原文明确``参数仅由训练数据估计，验证样本只用于调参''，
本文遵循此规定：验证集只用于早停与振幅校准。

\subsection{超参数}

超参数取自原文 Internet Appendix 表 A.5，而非凭经验设定（表 \ref{tab:hp}）。
本文亦测试了网格另一端的取值，主结果差异在 $0.05$ 个百分点以内，说明模型对此不敏感。

\begin{table}[htbp]
\centering
\caption{NN 系列超参数（源自原文 Internet Appendix 表 A.5）}
\label{tab:hp}
\small
\begin{tabular}{lll}
\toprule
超参数 & 原文网格 & 本文取值 \\
\midrule
$L_1$ 惩罚 $\lambda_1$ & $(10^{-5},\,10^{-3})$ & $10^{-5}$ \\
学习率 & $\{0.001,\ 0.01\}$ & $0.001$ \\
批大小 / 最大轮数 / 早停耐心 & $10{,}000$ / $100$ / $5$ & 同左 \\
Adam 参数 & 默认 & 默认 \\
集成规模 & $10$ 个随机种子 & $3$ 个随机种子 \\
批量归一化位置 & ReLU 之后 & ReLU 之后 \\
\bottomrule
\end{tabular}
\end{table}

\subsection{两项关键实现细节}

\textbf{早停准则。}以验证集均方误差早停在本问题上完全失效：年度收益的跨月均值差异极大，
使得每轮验证均方误差在四位小数上几乎相同，早停恒在第 1 轮触发，网络停留在初始化状态。
本文改用\textbf{月内横截面信息系数}作为准则——先对每月内的预测与目标去均值再计算合并相关，
从而剔除跨月水平差异。该改动使最优轮数从 1 提升到 8--20 轮，
验证集信息系数由约 $0.05$ 升至 $0.11$--$0.13$。

\textbf{预测振幅。}网络原始输出的量纲随输入分布漂移，把验证集拟合的缩放系数直接套用到
测试年会出现系统性错配（2008 年预测标准差为最优值的三倍，单年即损失约 $1$ 个百分点）。
本文改为逐月标准化后再施加目标离散度：
\begin{equation}
\hat y_{i,t} = \bar y_{\mathrm{train}} + \frac{\hat f_{i,t} - \overline{\hat f}_t}{\mathrm{sd}_t(\hat f)}
\cdot \left(\rho_{\mathrm{val}}\cdot \sigma_{y,\mathrm{val}}\right),
\end{equation}
即量化投资中标准的``信息系数乘以目标波动率''定额法则，仅使用验证集信息。

\subsection{评价指标}

样本外 $R^2$ 按原文式（19）：
\begin{equation}
\R = 1 - \frac{\sum_{(i,t)\in\mathcal{T}}\left(y_{i,t} - \hat y_{i,t}\right)^2}
{\sum_{(i,t)\in\mathcal{T}} y_{i,t}^2},
\end{equation}
分母为\textbf{未去均值}的收益平方和，即基准是``预测为零''而非``预测样本均值''。
若对分母去均值，$R^2$ 的含义会变为``相对历史均值的改进''，数值将系统性偏低。
本文严格采用未去均值口径，并按市值分组报告大市值组（每月前 1{,}000）与小市值组（后 1{,}000）。

\section{结果}\label{sec:results}

\subsection{月度口径：可复现性的直接证据}

月度目标是判断整条链路是否正确的最直接证据，因为原文表 1 报告了同一口径的数字，
且存在独立第三方复现可供对照。测试期样本 2{,}491{,}790 条观测。

\begin{table}[htbp]
\centering
\caption{月度样本外 $R^2$（论文表 1 口径，1987--2016）}
\label{tab:monthly}
\small
\begin{tabular}{lccc}
\toprule
 & 全样本 & 大市值组 & 小市值组 \\
\midrule
本文 NN3 & $+0.47\%$ & $+0.47\%$ & $+0.80\%$ \\
GKX (2020) NN3 & $0.40\%$ & $0.70\%$ & $0.45\%$ \\
第三方预注册复现 NN3 & $0.454\%$ & $0.294\%$ & $0.844\%$ \\
\bottomrule
\end{tabular}
\end{table}

三者处于同一水平，差异在统计不确定性之内（第三方报告的 $95\%$ 置信区间为
$[0.165\%, 0.733\%]$）。更值得注意的是结构上的一致性：\textbf{大市值组均低于原文，
小市值组均高于原文}。这种``同方向、同幅度''的偏离指向数据版本或股票池的细微差异，
而非实现错误。

月度结果覆盖了从设计矩阵读取、滚动切分、小批量训练、早停、预测、振幅定额到 $R^2$ 计算的
全部环节。任何一个环节存在缺陷，都不可能在月度口径上恰好落在两个独立来源之间。
\textbf{这是本文最重要的单项证据。}

\subsection{年度口径：主结果}

\begin{table}[htbp]
\centering
\caption{年度样本外 $R^2$（论文表 2 口径，1987--2016，$n=2{,}235{,}775$）}
\label{tab:annual}
\small
\begin{tabular}{lcccc}
\toprule
模型 & 全样本 & 大市值组 & 小市值组 & 正收益年份 \\
\midrule
\multicolumn{5}{l}{\emph{缩尾目标（$0.5\%/99.5\%$），$R^2$ 亦按缩尾后目标计算}} \\
NN3（3 种子集成，振幅校准） & $\mathbf{+2.71\%}$ & $+4.29\%$ & $+3.88\%$ & 21/30 \\
NN3（3 种子集成） & $+2.61\%$ & $+3.69\%$ & $+3.91\%$ & 21/30 \\
NN3（单种子） & $+2.50\%$ & $+3.44\%$ & $+3.81\%$ & 21/30 \\
常数均值基准 & $+1.97\%$ & $+4.58\%$ & $+2.70\%$ & -- \\
三特征线性基准 & $+1.94\%$ & $+2.84\%$ & $+3.17\%$ & -- \\
岭回归（920 维，$\lambda=0.1$） & $+1.66\%$ & -- & -- & -- \\
\midrule
\multicolumn{5}{l}{\emph{原始目标（论文口径）}} \\
NN3（3 种子集成，缩尾训练） & $\mathbf{+1.61\%}$ & -- & -- & -- \\
NN3（原始目标训练） & $+1.32\%$ & $+0.84\%$ & $+2.62\%$ & 19/30 \\
十特征线性基准 & $+1.04\%$ & $+1.74\%$ & $+1.82\%$ & -- \\
\midrule
GKX (2020) NN3 & $+3.40\%$ & $+4.73\%$ & $+5.17\%$ & -- \\
\bottomrule
\end{tabular}
\end{table}

表 \ref{tab:annual} 有两点与原文一致，一点不一致。

\textbf{一致之一：神经网络优于线性基准。}原始口径下 NN3 的 $+1.32\%$ 高于十特征线性基准的
$+1.04\%$；缩尾口径下 NN3 的 $+2.61\%$ 高于岭回归的 $+1.66\%$。

\textbf{一致之二：大市值组优于全样本。}缩尾口径下全样本 $+2.71\%$、大市值组 $+4.29\%$，
与原文排序一致。

\textbf{不一致：绝对水平。}缩尾口径低约 $0.7$ 个百分点，原始口径低约 $1.8$ 个百分点。
后文两节分别从数据边界与文献现状分析这一差距。

一个必须同时报告的对照是\textbf{常数均值基准}。因分母未去均值，
``预测历史均值''本身即可获得 $+1.97\%$（缩尾）或 $+0.75\%$（原始）。
换言之，本文 NN3 的增益中约 $1.7$ 个百分点来自均值成分，横截面预测的净贡献只有
$0.5$ 至 $0.9$ 个百分点。若报告年度 $R^2$ 时不列常数基准，数值是不可解释的——
一个毫无横截面预测能力的模型也能拿到接近 $2\%$ 的``成绩''。

\subsection{差距分析：为什么不是实现缺陷}\label{sec:gap}

本文用三项互相独立的证据回答``差距从哪来''。

\textbf{证据一：管线能恢复已知信号。}构造 $y = 0.030\,c_1 + 0.015\,c_2 + \varepsilon$
（$\varepsilon$ 为标准差 $0.06$ 的噪声，$c_1,c_2$ 为两个特征），
其可达 $R^2$ 可精确计算为 $+10.68\%$。用完全相同的管线（同样的特征、
同样的切分、同样的网络、同样的定额与评价函数）跑一遍，恢复 **$+10.60\%$**，
恢复率 **$0.99$**。若训练、早停、定额或 $R^2$ 计算任一环节存在缺陷，
不可能恰好命中理论值。

\textbf{证据二：年度模型已超出月度信号的上界。}用月度模型检验信号持续性：
月度预测与未来第 $k$ 月收益的横截面秩相关在 $k=1$ 时为 $0.037$，
此后稳定在 $0.024$--$0.031$ 直至 $k=12$，即信号持续但不强。
把这一持久信号累加到年度，可支撑的年度信息系数约为 $\sqrt{12}\times 0.025 \approx 0.087$，
对应 $R^2$ 约 $+0.84\%$。直接计算月度预测与 12 个月年度收益的秩相关得 $0.0916$，与之吻合。
本文年度模型达到 $+1.32\%$（信息系数 $0.11$），\textbf{已高于该上界}，
说明它确实学到了月度模型无法捕捉的长期成分。

\textbf{证据三：论文数值要求的信息系数超出了数据的支撑能力。}
$+3.40\%$ 对应信息系数 $\sqrt{0.034}\approx 0.185$，约为上述上界的两倍。
在这一约束下，差距不能仅靠改进训练过程或调参消除。

作为补充，本文系统排查了多种目标构造与正则化设定（表 \ref{tab:ablate}）。
无论采用哪一种，年度 $R^2$ 都稳定落在 $1.3\%$--$2.7\%$ 区间。

\begin{table}[htbp]
\centering
\caption{年度口径的方案排查}
\label{tab:ablate}
\small
\begin{tabular}{p{6.4cm}p{3.4cm}p{4.6cm}}
\toprule
方案 & 年度 $R^2$ & 说明 \\
\midrule
缩尾目标 $0.5\%$ + 3 种子集成 & $+2.71\%$ & 最优（缩尾口径） \\
原始目标直接训练 & $+1.32\%$ & 与论文口径完全对齐 \\
排名目标（横截面秩为训练目标） & $+0.03\%$ & 验证集 IC 升至 $0.21$，但对原始收益解释力弱 \\
月度预测自下而上加总 & $-0.78\%$ & 持续性假设不成立 \\
$168$ 维（94 特征 + 74 行业） & $+2.37\%$ & 缩尾口径 \\
$920$ 维（原文完整设计） & 验证集 IC $0.015$--$0.042$ & 显著劣于 94 维 \\
评估期缩尾 $0\to10\%$ 扫描 & 峰值 $2.71\%$ @ $0.5\%$ & 尾部不是瓶颈 \\
$L_1$ 惩罚 $10^{-3}$ & $+2.44\%$ & 与 $10^{-5}$ 无实质差异 \\
\bottomrule
\end{tabular}
\end{table}

其中两点值得说明。

\textbf{缩尾的作用与其边界。}年度目标尾部极厚：原始口径下 $E[y^2] = 0.836$，
仅前 10 个观测就占平方和的 $22\%$。这些极端值经核查是真实的——
例如某只股票 1999 年 4 月月度收益为 $+452\%$，价格由 \$4.75 升至 \$26.25，
此后升至 \$346，价格路径与收益率完全自洽，属互联网泡沫时期的小盘股，并非复权错误。
$0.5\%/99.5\%$ 缩尾使 $E[y^2]$ 降至 $0.369$，$R^2$ 升至 $+2.71\%$；
但缩尾并非越强越好，推至 $2\%$ 后反而降至 $+2.35\%$，说明模型确实捕捉到了部分大波动信息。

\textbf{原文的 OLS 基准隐含正则化。}本文用解析对照确认了这一点：
把岭回归强度 $\lambda$ 从 $10^{-10}$ 扫到 $10$，无正则的最小二乘解预测标准差达 $521$
（目标标准差 $0.52$），$R^2$ 为 $-1.2\times10^{8}\%$；
在 $\lambda \approx 5\times10^{-4}$ 时 $R^2$ 恰为 $-30.6\%$，与原文的 $-34.86\%$ 接近。
\textbf{裸最小二乘在公开数据上不可复现。}

\subsection{经济价值}

按原文做法构造十分位组合，各组实现年度超额收益见表 \ref{tab:decile}。

\begin{table}[htbp]
\centering
\caption{按 NN3 预测分组的实现年度超额收益（等权，\%）}
\label{tab:decile}
\small
\begin{tabular}{lcccccccccc}
\toprule
分位 & D1 & D2 & D3 & D4 & D5 & D6 & D7 & D8 & D9 & D10 \\
\midrule
收益 & $-3.09$ & $2.86$ & $5.26$ & $7.09$ & $8.35$ & $8.72$ & $9.26$ & $10.71$ & $12.87$ & $17.42$ \\
\bottomrule
\end{tabular}
\end{table}

十分位严格单调，仅最低组为负，多空价差为每年 $+20.51\%$。
以重叠年度收益计算的月度多空序列，年化 Sharpe 比率为 $1.16$；
论文为 $1.35$，第三方预注册复现为 $1.43$（$95\%$ 置信区间 $[1.06,1.82]$）。
\textbf{经济价值复现到 $86\%$，好于统计指标的 $80\%$}，
与``排序质量比拟合精度更容易迁移''的直觉一致。

\section{与公开复现文献的对照}\label{sec:lit}

本文调查了 GitHub 上六个公开的 GKX 复现项目（表 \ref{tab:lit}）。
目的不是寻找``标准答案''，而是判断本文的困难是否属于个别现象。

\begin{table}[htbp]
\centering
\caption{公开复现项目的口径与结果}
\label{tab:lit}
\small
\begin{tabular}{p{4.4cm}p{1.9cm}p{3.3cm}p{4.4cm}}
\toprule
项目 & 口径 & NN3 结果 & 备注 \\
\midrule
PatrickDamouni (2026) & 月度 & $0.454\%$ & 预注册、单次窥视、含密封留出集 \\
dinobao / HKUST & 月度 & $0.13\%$ & 课程项目 \\
RANN9 & 月度 & -- & notebook 内无年度目标代码 \\
saeid-torabi & 月度 & -- & ``年度 $R^2$''实为按年分组 \\
Tidy Finance (Voigt) & 月度 & -- & 仅拟合首个测试年 \\
Kratos0806 & 月度 & -- & 随机森林 + 组合绩效 \\
\midrule
本文 & 月度 & $+0.47\%$ & 与第三方一致 \\
本文 & 年度 & $+2.71\%$ / $+1.61\%$ & 尚无第三方报告 \\
\bottomrule
\end{tabular}
\end{table}

三点观察。

\textbf{其一，六个项目全部针对月度口径，无一报告年度数值。}
部分项目的文件名含 ``annual $R^2$''，实为把月度预测按年份分组计算，
与论文表 2 的 12 个月前瞻目标并非同一概念。这提示年度口径的复现难度显著高于月度。

\textbf{其二，最严谨的第三方也未能复现论文的线性基准。}PatrickDamouni 在结论中明确写道，
其 OLS 全变量回归的 $R^2$ 为 $-170\%$，而原文为 $-3.5\%$，并列为``未解决的设定差异''。
这与本文在年度口径上的困难（裸最小二乘发散至 $10^8$ 量级，
须加 $\lambda\approx5\times10^{-4}$ 才得到与原文接近的 $-30.6\%$）是同一现象。
两个独立团队在不同口径上遇到同类问题，说明这是公开数据与论文之间的一项共性差异。

\textbf{其三，第三方观察到的结构性现象与本文一致}：模型排序发生改变
（第三方的 NN1 优于 NN3，而原文认为 NN3 最优）；单一种子不足以超越常数基准，必须依赖集成；
小市值组 $R^2$ 普遍高于原文，大市值组普遍低于原文。这些一致性进一步支持本文月度结果的可信度。

\section{结论}\label{sec:conclusion}

\textbf{月度口径可以复现。}本文 NN3 取得 $+0.47\%$，论文为 $0.40\%$，
独立第三方预注册复现为 $0.454\%$，三者互相印证。这一结果覆盖了从数据读取到
$R^2$ 计算的全部环节，是管线正确性的直接证据。

\textbf{年度口径未能复现，但差距不是实现缺陷。}本文最优结果为缩尾口径 $+2.71\%$、
原始口径 $+1.61\%$，论文为 $+3.40\%$。三项证据表明差距不可归因于代码：
合成信号注入测试中管线恢复了理论上限的 $99\%$；由月度信号推出的年度可预测性上界约为
$+0.84\%$，而本文年度模型已达 $+1.32\%$，超出该上界；
论文数值要求的信息系数 $0.185$ 约为月度信号可持续支撑水平的两倍。

\textbf{方向性结论全部复现。}非线性优于线性基准、大市值组优于全样本、
动量与流动性及波动率是最重要的信号、十分位组合严格单调、经济价值显著
（多空 Sharpe $1.16$ 对论文 $1.35$），均与原文一致。统计指标复现到 $80\%$，
经济价值复现到 $86\%$。

\textbf{公开文献中尚无年度复现的先例。}六个公开项目全部针对月度口径，
其中最严谨的一项亦未能复现论文的线性基准，与本文困难同源。
据此本文认为，年度 $R^2$ 的未达更可能反映公开数据条件下的可复现性边界，
而非某一实现的缺陷。

\textbf{局限。}第一，本文未能获得原文表 2 的计算细节，若其年度目标构造与本文不同
（例如采用收益加总而非复利），数值对比将失去直接意义。第二，主结果使用 94 维特征
而非原文的 920 维，尽管对照显示 920 维在本框架下表现更差，该偏离仍需说明。
第三，多空组合的 Sharpe 比率基于重叠年度收益计算，与按月重构的组合并非完全可比。

\appendix

\section{代码与复现材料}\label{app:code}

全部代码、配置与本文所用验证脚本已发布于：

\begin{center}
\url{\repourl}
\end{center}

仓库中包含完整的滚动重估实现、13 个验证与诊断脚本、
一份可直接编译的 \LaTeX{} 源文件，以及本文的 PDF。
因篇幅所限，报告不再附印全部源码，此处仅列出文件职责（表 \ref{tab:files}）。

\begin{table}[htbp]
\centering
\caption{仓库文件职责}
\label{tab:files}
\small
\begin{tabular}{p{4.6cm}p{9.4cm}}
\toprule
文件 & 职责 \\
\midrule
\code{config.json} & 全部可调参数：数据路径、目标口径、特征维度、网络与训练超参数 \\
\code{src/prep.py} & 读取特征文件、合并 CRSP 收益、构造 $y^{(1)}$ 与 $y^{(12)}$、缓存面板 \\
\code{src/features.py} & 横截面排序、中位数填补、宏观交互、行业虚拟变量、按年写出设计矩阵 \\
\code{src/nn.py} & 网络定义、Adam 训练、早停、振幅校准 \\
\code{src/train.py} & 滚动重估、目标构造、$R^2$ 计算、线性基准 \\
\code{run.py} / \code{probe.py} & 完整运行与快速试跑 \\
\code{ensemble.py} & 多种子预测平均与振幅扫描 \\
\code{audit\_*.py} & 数据与代码正确性验证（见附录 \ref{app:audit}） \\
\code{probe\_*.py} & 诊断性探针：持续性、缩尾灵敏度、自下而上聚合、OLS 正则对照 \\
\bottomrule
\end{tabular}
\end{table}

运行方式：先 \code{python run.py prep} 生成缓存与设计矩阵，
再 \code{python run.py nn --seeds <种子> --tag <标签>} 完成一次完整滚动重估，
结果写入 \code{outputs/<tag>\_result.json} 与 \code{outputs/<tag>\_preds.npz}。
配置中 \code{data.target} 可取 \code{annual}、\code{monthly} 或 \code{annual\_rank}，
\code{model.n\_features} 可取 \code{94}、\code{"chars+ind"} 或留空（920 维）。

\section{与原文的偏离清单}\label{app:deviation}

\begin{table}[htbp]
\centering
\small
\begin{tabular}{p{0.5cm}p{4.4cm}p{4.4cm}p{4.4cm}}
\toprule
序号 & 偏离 & 原因 & 影响评估 \\
\midrule
1 & 特征维度取 94 而非 920 & 课程要求；920 维在本框架下验证集 IC 仅 $0.015$--$0.042$ & 主结果口径；对照见 \ref{tab:ablate} \\
2 & 集成规模 3 而非 10 & 计算资源限制 & 单种子到 3 种子提升 $0.11$ 个百分点 \\
3 & 年度目标做 $0.5\%/99.5\%$ 缩尾 & 原始目标尾部极厚，L2 目标被少数观测主导 & 两种口径均报告 \\
4 & 早停用月内信息系数而非验证集 MSE & 后者在本问题上不可用 & 使模型能够收敛 \\
5 & 振幅按信息系数乘目标波动率定额 & 网络输出量纲随输入分布漂移 & 2008 年单年损失减少约 $1$ 个百分点 \\
6 & 批量归一化置于 ReLU 之后 & 依原文 Internet Appendix 算法 7 & 与原文一致 \\
7 & 测试期止于 2016 年 12 月 & 特征文件止于该月 & 与原文一致 \\
8 & 使用 Fama--French 无风险利率 & 与 CRSP 国库券利率有细微差异 & 影响可忽略 \\
\bottomrule
\end{tabular}
\end{table}

其中第 3 项最需说明。原文未提及对收益做任何缩尾或剔除，
本文亦以原始目标训练过（$+1.32\%$），但缩尾训练配合原始口径评价可得 $+1.61\%$，
说明适度缩尾起到了正则化作用。本文因此同时报告两种口径，而不隐藏任何一种。

\section{验证性审计}\label{app:audit}

本节列出用于自证的脚本。它们不产生结果，只用于排除``结果不可信''的可能（表 \ref{tab:auditscripts}）。

\begin{table}[htbp]
\centering
\caption{验证与诊断脚本}
\label{tab:auditscripts}
\small
\begin{tabular}{p{4.2cm}p{9.8cm}}
\toprule
脚本 & 检验内容 \\
\midrule
\code{audit\_synth.py} & 合成信号注入：构造已知信号的目标，检验管线能否按理论值恢复。实测 $+10.60\%$ 对理论上限 $+10.68\%$，恢复率 $0.99$ \\
\code{audit\_prep.py} & 从原始 CSV 独立重算某月横截面排名，与缓存设计矩阵比对 \\
\code{audit\_md\_rf.py} & 按教师给出的公式校验 \code{md} 的解码（1{,}156 个月份零分歧），并将所用无风险利率与 Fama--French 官方文件逐月比对（718 个月，最大差 $10^{-16}$） \\
\code{audit\_y12.py} & 手工计算 12 个月连乘，与缓存的年度目标比对 \\
\code{audit\_target.py} & 用三个特征的信息系数判定年度目标窗口起点，排除错位一月 \\
\code{audit\_fit.py} & 报告训练集、验证集、测试集三档信息系数，区分过拟合与欠拟合 \\
\code{audit\_counts.py} & 逐月统计三个数据源的股票数，解释与原文 $6{,}200$ 的差异 \\
\code{probe\_ols\_ridge.py} & 扫描岭回归强度，定位原文 OLS 基准对应的 $\lambda$ \\
\code{probe\_persist.py} & 测量月度信号在 $k=1\ldots12$ 个预测期上的衰减 \\
\code{probe\_clipsens.py} & 评估期缩尾强度的灵敏度扫描 \\
\code{probe\_bottomup.py} & 由月度预测自下而上构造年度预测并评价 \\
\bottomrule
\end{tabular}
\end{table}

\end{document}
